% 编译方式: xelatex SL_fractional_left_definite.tex
\documentclass[UTF8]{ctexart}
\usepackage{amsmath, amssymb, amsthm}
\usepackage[margin=2.5cm]{geometry}
\usepackage[hyphens]{url}
\usepackage[hidelinks]{hyperref}
\usepackage{booktabs}
\emergencystretch 3em

\newtheorem{theorem}{定理}[section]
\newtheorem{proposition}[theorem]{命题}
\newtheorem{lemma}[theorem]{引理}
\newtheorem{corollary}[theorem]{推论}
\newtheorem{remark}[theorem]{注}
\newtheorem{definition}[theorem]{定义}

\title{稀疏多项式族的分数阶稠密性: $0\le s\le3$ 与高阶推广勘误}
\author{研究证明文档 (项目: BVE research)}
\date{2026-08-05; 2026-09-20 revision}

\begin{document}
\maketitle

\begin{abstract}
2026-09-20 修订. 旧版声称同一稀疏多项式族在所有实数阶 $H^s$ 中完备,
该全范围结论已撤回: 对 $s\ge0$, $p_4=x^4-2x^2\in H^s$ 当且仅当 $s<7/2$.
本版仅保留由 H3 主定理、谱截断和连续嵌入严格推出的 $0\le s\le3$
结论, 它覆盖原分数窗口 $3/2\le s<2$. $3<s<7/2$ 的同族稠密性本次不判定.
原文与历史数值证据按修订前字节保存在本轮审计工件中; 它们不再作为高阶幂域证明.
\end{abstract}

\tableofcontents

\section{对象与可用输入}
固定 $c>0$, 令 $K_cf=-f''+cf$ 为 $L^2(-1,1)$ 中的移位 Krein Laplacian,
\[
 D(K_c)=\{f\in H^2_{\rm Sob}(-1,1):
 f'(1)=f'(-1)=(f(1)-f(-1))/2\}.
\]
这里 $H^2_{\rm Sob}$ 表示通常 Sobolev 空间, 避免与左定记号混淆.
对实数 $s\ge0$ 定义
\[
 H^s=D(K_c^{s/2}),\qquad \|f\|_{H^s}=\|K_c^{s/2}f\|_2.
\]
自伴性与 $K_c\ge cI$ 为既有左定设定. 完整正交谱系及其 $L^2$ 范数平方为
\[
\begin{array}{c|c|c}
 \phi & \lambda & \|\phi\|_2^2\\ \hline
 1 & c & 2\\
 x & c & 2/3\\
 \cos(n\pi x),\ n\ge1 & (n\pi)^2+c & 1\\
 \sin(\mu_nx),\ n\ge1 & \mu_n^2+c & \mu_n^2/(1+\mu_n^2)
\end{array}
\]
其中 $n\pi<\mu_n<(n+1/2)\pi$, $\tan\mu_n=\mu_n$.
特别地 $c$ 的重数为二. 偶子空间在 $[0,1]$ 上给出 Neumann--Neumann
问题; 奇子空间给出 $f(0)=0$, $f'(1)=f(1)$ 的问题, 后者的零频模态为 $x$.
两者都是正则分离边界问题, 其完备正交系合起来即上表. 直接积分还给出
\[
 \int_{-1}^1x\sin(\mu_nx)\,dx
 =\frac{2(\sin\mu_n-\mu_n\cos\mu_n)}{\mu_n^2}=0,\qquad
 \|\sin(\mu_nx)\|_2^2=1-\frac{\sin(2\mu_n)}{2\mu_n}
 =\frac{\mu_n^2}{1+\mu_n^2}.
\]
把上表归一化并枚举为 $\{e_j\}$, 对 $s\ge0$ 有真正的幂域判据
\[
 f\in H^s\quad\Longleftrightarrow\quad
 \sum_j\lambda_j^s|\langle f,e_j\rangle|^2<\infty,
 \qquad \|f\|_{H^s}^2=\sum_j\lambda_j^s|\langle f,e_j\rangle|^2.
\]
若使用未归一化的 $\phi_j$, 每项必须除以 $\|\phi_j\|_2^2$.

对 $s<0$, 定义 $H^s$ 为 $L^2$ 在范数 $\|K_c^{s/2}f\|_2$ 下的完备化,
即加权序列空间 $\{(a_j):\sum_j\lambda_j^s|a_j|^2<\infty\}$.
虽然有界算子 $K_c^{s/2}$ 在 $L^2$ 上处处有定义, 这个较弱范数不与
$L^2$ 范数等价: $\|\cos(n\pi x)\|_s=((n\pi)^2+c)^{s/2}\to0$,
而其 $L^2$ 范数恒为一. 事实上对任意 $s<0$, 取余弦模态系数
$a_n=n^{-1/2}$, 则 $\sum\lambda_n^s/n<\infty$ 而 $\sum1/n=\infty$;
其有限和在弱范数下 Cauchy, 在 $L^2$ 中没有该范数的极限.
这同时说明必须作完备化. 在这个完整标度上, 系数映射
\[
 \mathcal K_c:(a_j)\longmapsto(\lambda_ja_j),\qquad
 \mathcal K_c:H^s\longrightarrow H^{s-2}
\]
对每个实数 $s$ 是等距满射, 逆映射逐项乘以 $\lambda_j^{-1}$.
它在 $D(K_c)$ 上与原算子相同, 在其它阶上表示标度延拓, 不能默认
对所有元素都按经典二阶导数作用. 下述低阶稠密性证明仍不依赖负阶空间.

稀疏族为
\[
 p_0=1,\quad p_1=x,\quad
 p_{2m}=x^{2m}-\frac m{m-1}x^{2m-2},\quad
 p_{2m+1}=x^{2m+1}-\frac m{m-1}x^{2m-1},\quad m\ge2.
\]
每个 $p_n$ 满足边界条件, 且 $K_cp_n$ 为光滑多项式, 属于
$D(K_c^{1/2})=H^1_{\rm Sob}(-1,1)$, 所以 $p_n\in H^3$.
本页使用修订后的 \texttt{SL\_h3\_completeness\_proof.tex} 主定理:
$\operatorname{span}\{p_n\}$ 在 $H^3$ 中稠密. 该输入仍由 H1 矩跳跃主链证明,
不依赖已撤回的任意阶推广.

\section{谱截断传输的严格证明}
\begin{theorem}[低阶和分数窗口]\label{thm:main}
对每个 $c>0$ 和实数 $0\le s\le3$, $\{p_n\}$ 在 $H^s$ 中稠密.
\end{theorem}
\begin{proof}
设 $E$ 为 $K_c$ 的谱测度. 对 $g\in H^3$, 因谱包含于 $[c,\infty)$,
\[
 \|g\|_{H^s}^2=\int\lambda^s\,d\|E_\lambda g\|^2
 \le c^{s-3}\int\lambda^3\,d\|E_\lambda g\|^2
 =c^{s-3}\|g\|_{H^3}^2.
\]
取 $f\in H^s$ 及 $\varepsilon>0$. 截断 $f_N=E([c,N])f$ 属于 $H^3$,
因为谱支集有界. 由积分尾部收敛,
$\|f-f_N\|_{H^s}^2=\int_{(N,\infty)}\lambda^s\,d\|E_\lambda f\|^2\to0$.
先取 $N$ 使此误差范数小于 $\varepsilon/2$, 再由 H3 主定理取有限线性组合
$p\in\operatorname{span}\{p_n\}$ 使
$\|f_N-p\|_{H^3}<\varepsilon/(2c^{(s-3)/2})$.
三角不等式和连续嵌入给出 $\|f-p\|_{H^s}<\varepsilon$. \qed
\end{proof}

\section{反例与撤回的证明机制}
\begin{proposition}[精确成员阈值]\label{prop:thresholds}
对每个 $c>0$, $t,s\ge0$,
\[
 x^2\in H^t\iff t<3/2,\qquad p_4\in H^s\iff s<7/2.
\]
因此 $s\ge7/2$ 时原稀疏族甚至不是 $H^s$ 的子集.
\end{proposition}
\begin{proof}
对 $a=n\pi$, 两次分部积分给出
\[
 \int_{-1}^1x^2\cos(ax)\,dx=\frac{4(-1)^n}{a^2},\qquad
 \int_{-1}^1x^4\cos(ax)\,dx
 =\frac{8(-1)^n}{a^2}-\frac{48(-1)^n}{a^4}.
\]
所以 $p_4$ 的 $a^{-2}$ 项恰好相消. 与未归一化常数模态的内积分别为 $2/3$ 与
$-14/15$, 所有奇模态的投影为零. 允许级数值为 $+\infty$, 得到
\begin{align*}
 \|x^2\|_t^2&=\frac{2c^t}{9}
 +16\sum_{n\ge1}\frac{((n\pi)^2+c)^t}{(n\pi)^4},\\
 \|p_4\|_s^2&=\frac{98c^s}{225}
 +2304\sum_{n\ge1}\frac{((n\pi)^2+c)^s}{(n\pi)^8}.
\end{align*}
尾项分别与 $n^{2t-4}$ 和 $n^{2s-8}$ 等价. 正项级数判别给出上述
严格阈值, 两个端点均为调和级数型发散. 这里 $+\infty$ 只表示不属于
相应空间, 不是空间元素的有限范数. \qed
\end{proof}

保留第一轮的独立算子域反例作为一致性检查:
$p_4\in D(K_c)$, 但
\[
 K_cp_4=cx^4-(2c+12)x^2+4,\qquad (K_cp_4)'(1)=-24,\quad
 (K_cp_4)'(-1)=24.
\]
偶函数 $K_cp_4$ 的边界值差为零, 两端导数却不为零, 所以不属于 $D(K_c)$.
按算子乘积域定义 $p_4\notin D(K_c^2)$, 与精确阈值相符.

旧版的 $N_k=(w,x^k)_t$ 任意阶取矩不成立: $x^2\notin D(K_c)$,
所以高阶内积可能未定义. 同样, $(c-D^2)^r$ 的形式多项式展开不证明
真正算子幂存在. 稠密性传输只可应用于已经属于源空间的族.
尤其从 $t=3/2$ 起, 连 $x^2$ 的 $H^t$ 矩都不能这样定义.
负阶标度虽可按上文修正, 却不消除正阶的成员障碍.
历史脚本的投影残差或形式范数增长不是幂域成员关系的证书.
$p_4$ 的成员判据也不是 $3<s<7/2$ 时整个族的稠密性证明.

\section{仍然有效的增长引理}
\begin{lemma}[有显式系数条件的钉住解]\label{lem:growth}
设 $c_0>0$, $u_0=0$, $u_1=1$, 且
$c_0u_j=A_ju_{j-1}-B_ju_{j-2}$ ($j\ge2$).
若 $B_j\ge0$ 且 $A_j-B_j\ge c_0$ 对每个 $j\ge2$ 成立, 则
$u_j>0$ ($j\ge1$), $u_j\ge u_{j-1}$, 并且
\[
 u_j\ge\prod_{k=2}^j\frac{A_k-B_k}{c_0}.
\]
若还满足 $A_j-B_j\ge4j+c_0$, 则
$u_j\ge(4/c_0)^{j-1}j!$, 因而增长快于每个固定次数的多项式.
\end{lemma}
\begin{proof}
从 $0=u_0<u_1=1$ 出发, 归纳使用
\[
 c_0u_j=(A_j-B_j)u_{j-1}+B_j(u_{j-1}-u_{j-2})
 \ge(A_j-B_j)u_{j-1}\ge c_0u_{j-1}.
\]
逐项连乘得第一下界; 再用 $(A_j-B_j)/c_0\ge4j/c_0$ 得第二下界.
仅有前两个条件不保证快于多项式的增长: $A_j=c_0$, $B_j=0$ 时
$u_j=1$ ($j\ge1$). \qed
\end{proof}

本问题的偶、奇递推系数分别为
\[
 A_j=2j(2j-1)+\frac{cj}{j-1},\quad B_j=2j(2j-3),\qquad
 A'_j=2j(2j+1)+\frac{cj}{j-1},\quad B'_j=2j(2j-1).
\]
两者均满足 $A_j-B_j=A'_j-B'_j=4j+cj/(j-1)\ge4j+c$.
引理由直接归纳成立, 不依赖任何衰减根断言. 更确切地, 若某个递推解的
比值 $r_j=u_j/u_{j-1}$ 最终有定义并收敛到有限非零 $L$, 则
$cr_j=A_j-B_j/r_{j-1}$ 除以 $j^2$ 后给出 $0=4-4/L$, 故只能有 $L=1$.
这不声称这样的解存在或决定其增长阶. 增长引理的应用仍须先确认矩内积
有定义及其增长上界; 此处保留的 H3 证明只在 $H^1$ 中取矩.

\section{数值工具的含义与限制}
\path{scripts/op10_verify_mechanism.py} 与
\path{scripts/op10_fractional_window.py} 中的有限谱和只表示谱投影的平方范数,
不自动表示原函数的有限范数. 对命题 \ref{prop:thresholds} 的发散区间,
程序应明确报告解析发散, 不以有限截断宣称通过. 在收敛区间可使用尾界:
若 $q-2t>1$, 则
\[
 \sum_{n>N}\frac{((n\pi)^2+c)^t}{(n\pi)^q}
 \le \pi^{2t-q}\left(1+\frac{c}{\pi^2(N+1)^2}\right)^{\max(t,0)}
 \frac{N^{2t-q+1}}{q-2t-1}.
\]
这是用单调因子上界和 $p$ 级数的积分尾界得到的解析估计; 普通浮点计算
该表达式仍不等于带向外舍入的区间证书.

为稳定计算第 $k$ 个正根, 写 $P_k=(k+1/2)\pi$, $\delta=P_k-\mu$,
在 $[0,\pi/2]$ 求
\[
 h_k(\delta)=\cos\delta-(P_k-\delta)\sin\delta=0.
\]
有 $h_k(0)=1$, $h_k(\pi/2)=-k\pi<0$,
$h_k'(\delta)=-(P_k-\delta)\cos\delta<0$ 于开区间, 因而严格括住唯一根.
程序检查端点符号、分支和缩放残差; 若浮点精度已不能分辨根与极点,
应显式失败, 而非把极点或固定偏移端点当根. 这些诊断、自检和有限传输
恒等式均不是独立验收或无限维稠密性的数值证明.

\begin{thebibliography}{9}
\bibitem{axioms} A. Jones, L. L. Littlejohn, A. Quintero Roba, \emph{Krein-Sobolev
	Orthogonal Polynomials II}, Axioms 14 (2025), 115.
	\url{https://doi.org/10.3390/axioms14020115}
\bibitem{lw1} L. L. Littlejohn, R. Wellman, \emph{A general left-definite theory
	for certain self-adjoint operators with applications to differential equations},
	J. Differential Equations 181 (2002), 280--339.
\bibitem{fg} C. Fischbacher, F. Gesztesy, P. Hagelstein, L. L. Littlejohn,
	\emph{Abstract left-definite theory: a model operator approach, examples,
	fractional Sobolev spaces, and interpolation theory}, arXiv:2408.01514 (2024).
	\url{https://arxiv.org/abs/2408.01514}
\bibitem{hs} 项目文档, \emph{左定空间 $H^s[-1,1]$ 中显式完备正交多项式系: 证明文档},
	2026-08-05.
\bibitem{dc} 项目文档, \emph{边界约束 Hilbert 空间中多项式稠密的一般准则},
	2026-08-05.
\end{thebibliography}

\end{document}
